% Standalone problem document, generated from the adjacent README.md.
% Compile with XeLaTeX. Mathematical content is maintained in Markdown.
\documentclass[11pt,a4paper]{article}
\usepackage[fontsize=10.5pt]{scrextend}
\usepackage[margin=22mm,headheight=15pt,headsep=9mm,footskip=11mm]{geometry}
\usepackage{fontspec}
\setmainfont{texgyrepagella}[Extension=.otf,UprightFont=*-regular,BoldFont=*-bold,ItalicFont=*-italic,BoldItalicFont=*-bolditalic]
\setsansfont{texgyreheros}[Extension=.otf,UprightFont=*-regular,BoldFont=*-bold,ItalicFont=*-italic,BoldItalicFont=*-bolditalic]
\setmonofont{texgyrecursor}[Extension=.otf,UprightFont=*-regular,BoldFont=*-bold,ItalicFont=*-italic,BoldItalicFont=*-bolditalic,Scale=MatchLowercase]
\usepackage{amsmath,amssymb}
\usepackage{unicode-math}
\setmathfont{texgyrepagella-math.otf}
\usepackage{xcolor,microtype,enumitem,needspace,fancyhdr}
\usepackage{bookmark}
\definecolor{ink}{HTML}{17344B}
\definecolor{accent}{HTML}{197D80}
\definecolor{muted}{HTML}{596773}
\hypersetup{colorlinks=true,linkcolor=accent,urlcolor=accent,pdftitle={TR-23 - Irreducibility
of asymptotic tensor-rank sublevel
varieties},pdfauthor={Open Problems in Numerical Linear Algebra}}
\urlstyle{same}
\setlength{\parindent}{0pt}
\setlength{\parskip}{5pt plus 1pt minus 1pt}
\linespread{1.04}
\setlength{\emergencystretch}{3em}
\widowpenalty=10000
\clubpenalty=10000
\displaywidowpenalty=10000
\allowdisplaybreaks[1]
\setlist{nosep,leftmargin=1.5em}
\providecommand{\tightlist}{\setlength{\itemsep}{0pt}\setlength{\parskip}{0pt}}
\providecommand{\pandocbounded}[1]{#1}
\pagestyle{fancy}
\fancyhf{}
\fancyhead[L]{\sffamily\footnotesize\color{muted}OPEN PROBLEMS IN NUMERICAL LINEAR ALGEBRA}
\fancyhead[R]{\sffamily\bfseries\footnotesize\color{accent}TR-23}
\fancyfoot[L]{\parbox[b]{0.9\textwidth}{\sffamily\fontsize{8}{10}\selectfont\color{muted}Editorial ratings; a bounded search cannot certify that a problem remains unsolved.\\Literature check: 2026-09-10}}
\fancyfoot[R]{\sffamily\footnotesize\color{muted}\thepage}
\renewcommand{\headrulewidth}{0.3pt}
\renewcommand{\headrule}{\hbox to\headwidth{\color{accent}\leaders\hrule height \headrulewidth\hfill}}
\makeatletter
\renewcommand{\section}{\Needspace{5\baselineskip}\@startsection{section}{1}{0pt}{12pt plus 2pt minus 2pt}{4pt}{\sffamily\bfseries\large\color{ink}}}
\renewcommand{\subsection}{\Needspace{4\baselineskip}\@startsection{subsection}{2}{0pt}{9pt plus 1pt minus 1pt}{3pt}{\sffamily\bfseries\normalsize\color{ink}}}
\makeatother
\setcounter{secnumdepth}{0}
\begin{document}
{\sffamily\small\color{muted}Tensor computations\par}
\vspace{3pt}
{\raggedright\hyphenpenalty=10000\exhyphenpenalty=10000\sffamily\bfseries\fontsize{21}{24}\selectfont\color{ink}Irreducibility
of asymptotic tensor-rank sublevel varieties\par}
\vspace{7pt}
{\sffamily\small\textbf{Difficulty:} extreme\quad\textbf{Importance:} interesting
to specialist\par}
{\sffamily\small\bfseries\color{accent}Status: Open\par}
\vspace{4pt}
{\color{accent}\hrule height 0.8pt}
\vspace{6pt}
\textbf{Rating rationale:} Controlling irreducibility for every
tensor-power cost sublevel is a foundational geometric question stronger
than finiteness of the attainable costs. Its immediate audience is
specialists in the algebraic geometry of tensor complexity.

\subsection{Problem statement}\label{problem-statement}

Let \(\mathbb F\) be an algebraically closed field, let \(k\geq3\), and
let \(d_1,\ldots,d_k\) be positive integers. Write
\(V=\mathbb F^{d_1}\otimes\cdots\otimes\mathbb F^{d_k}\). For
\(T\in V\), let \(R(T)\) be its tensor rank, and define its asymptotic
rank by \[
\widetilde R(T)=\lim_{m\to\infty}R(T^{\boxtimes m})^{1/m},
\] where \(\boxtimes\) groups corresponding modes of tensor powers; put
\(\widetilde R(0)=0\).

For every real \(a\geq0\), is the set \[
V_{\leq a}=\{T\in V:\widetilde R(T)\leq a\}
\] an irreducible algebraic variety?

The source proves that these sets are Zariski closed. The open assertion
is that none can be written as the union of two proper Zariski-closed
subsets. All fields, formats, and thresholds above are quantified;
\(0\in V_{\leq a}\) avoids the empty-set convention for irreducibility.

\subsection{Why it matters}\label{why-it-matters}

For matrices, bounded rank defines an irreducible determinantal variety.
The question asks whether a comparable geometric unity survives when
tensor rank is replaced by its asymptotic computational cost. A positive
answer would constrain the number of costs possible in a format and the
ways low-cost tensor families can degenerate.

This is stronger than merely asking for a finite set of asymptotic-rank
values. The source lists the two questions separately and records the
implication from irreducibility to discreteness; no converse is claimed.

\subsection{References and status
check}\label{references-and-status-check}

\begin{enumerate}
\def\labelenumi{\arabic{enumi}.}
\tightlist
\item
  M. Christandl, K. Hoeberechts, H. Nieuwboer, P. Vrana, and J. Zuiddam,
  \emph{Asymptotic tensor rank is characterized by polynomials},
  \href{https://arxiv.org/pdf/2411.15789v2}{arXiv:2411.15789v2}, Theorem
  1.2 and §5 second bullet, p.17.
\item
  Their \href{https://ir.cwi.nl/pub/36100/36100.pdf}{STOC 2025
  publication}, p.753, ``Geometric properties; irreducibility.''
\item
  J. Zuiddam,
  \href{https://staff.fnwi.uva.nl/j.zuiddam/talks/simons-14nov-2025.pdf}{November
  2025 research slides}, ``Open problems,'' question (2).
\end{enumerate}

\subsubsection{Status check ---
2026-09-10}\label{status-check-2026-09-10}

Checked the June 2026 arXiv revision, §5 second bullet, and the STOC
2025 irreducibility discussion, with targeted 2025--2026 searches. The
sources leave irreducibility open for tensors of order at least three.
Closedness is proved, but neither it nor the excluded matrix case
supplies a nontrivial resolved part of this universal irreducibility
target. Conditional implications from the full asymptotic-rank
conjecture are not unconditional solutions. No general proof or
counterexample was located.
\end{document}
